\chapter{Model Improvement on Real Data}
\label{chap:real_improvement}

\MicroTEDAClusPlus is indistinguishable from the worst model in six of the eight real datasets of \cref{chap:real_data}, \cref{tab:real_quality}. In this chapter, we repeat on these datasets the effort of \cref{chap:ModelImprovement}: we examine what the model does on them, evaluate changes to it, keep three, and confirm them on new runs and on datasets that took no part in their choice. The result is the final version of \MicroTEDAClusPlus, which we then compare with the benchmark models by the protocol of \cref{chap:real_data}, in quality and in cost. Where the two versions appear together, we call the model of \cref{chap:ModelImprovement} ``\MicroTEDAClusPlus of \cref{chap:ModelImprovement}'', and \modelname{MicroTEDAClus+ (Ch.~\ref{chap:ModelImprovement})} in the tables.

We diagnose the model and develop the changes on six of the eight real datasets, and leave \KDDCup and \Covertype out of both: they are the two longest streams, \cref{tab:real_datasets}, too costly for the many evaluations of the development, and in \Covertype the best model scores lower than in any other dataset, \cref{tab:real_quality}, which leaves little to separate the candidates. The two return in the confirmation, where they show whether the changes hold on data that took no part in their choice.

\section{Diagnosis}
\label{sec:ri_diagnosis}

\begin{table}[htbp]
\centering
\small
\caption[\modelname{MicroTEDAClus+ (Ch.~\ref{chap:ModelImprovement})} on the real datasets of the development]{\modelname{MicroTEDAClus+ (Ch.~\ref{chap:ModelImprovement})} on the real datasets of the development, in the runs of the real-data experiment: medians over the batches of the evaluation segment, then medians over 32 runs, except 31 in GAS. A mature microcluster has three samples or more.}
\adjustbox{max width=\textwidth, max totalheight=0.82\textheight}{%
\begin{tabular}{l c c c c c c }
\toprule
 & INS-A & INS-I & INS-R & GAS & PEN & SHU \\
\midrule
\multicolumn{7}{l}{\textit{Per batch}} \\
Share of the samples labelled as noise & 0.45 & 0.10 & 0.15 & 0.03 & 0.10 & 0.02 \\
Predicted clusters & 3.0 & 7.0 & 1.0 & 9.0 & 8.0 & 8.5 \\
Classes & 5.0 & 6.0 & 5.0 & 6.0 & 10.0 & 3.0 \\
Microclusters & 96.0 & 97.5 & 73.5 & 95.0 & 77.5 & 250.8 \\
Mature microclusters & 12.0 & 9.0 & 8.0 & 53.0 & 15.5 & 187.5 \\
Active microclusters & 11.0 & 8.2 & 4.0 & 28.0 & 12.0 & 83.8 \\
\bottomrule
\end{tabular}
}
\label{tab:ri_diagnosis}
\end{table}


\Cref{tab:ri_diagnosis} gives the state of \cref{chap:ModelImprovement}'s implementation \MicroTEDAClusPlus at every batch of the evaluation segments of the runs of \cref{chap:real_data}, on the six datasets of the development. In the INSECTS datasets and in \Pendigits, most of its microclusters never reach three samples, so few of them mature and fewer are active: the activation, the macroclusters and the prediction rest on a small share of the microclusters. In \GasDrift and \Shuttle, most microclusters mature. In \InsectsAbrupt the model labels nearly half of the samples as noise, and in \InsectsAbrupt and \InsectsReoccurring it predicts fewer clusters than there are classes, a single one in \InsectsReoccurring.

\section{Experimental Design}
\label{sec:ri_design}

The design is that of \cref{sec:mi_design}, with the changes that this section describes. We tune and evaluate every candidate by the protocol of \cref{sec:real_protocol}.

\subsection{Runs}
\label{subsec:ri_runs}

The runs of \cref{chap:real_data} are the runs 0 to 31. We develop the changes on the 32 runs that follow them, and every candidate uses, in each task, the seeds of the previous \MicroTEDAClusPlus version in that task, so that two candidates differ only by their changes. We confirm the kept changes on the runs 0 to 31 of the eight datasets, on which no candidate ran during the development. The INSECTS datasets and \Covertype keep the order of their files, \cref{subsec:real_order}, so every run of them processes the same stream, and their runs differ only in the seeds of the search.

\subsection{Decision Rule}
\label{subsec:ri_decision_rule}

We compare a candidate with its control as in \cref{subsec:mi_decision_rule}, on the Overall ARI of the six datasets and on the Steady-State ARI of the four whose penultimate batch does not hold a single class, \cref{tab:real_single_class}. Unlike in \cref{chap:ModelImprovement}, a candidate does not need to improve its control to be kept. We keep a candidate if

\begin{itemize}
    \item it improves its control;
    \item it is indistinguishable from its control, its change is analytically preferred, and its summary is not clearly larger;
    \item it is mixed, its gains outnumber its losses, and none of its losses, which we examine before the decision, comes from a defect of the change.
\end{itemize}

A kept candidate becomes the control of the candidates that follow it. As the cost of a candidate we report the size of its summary; we did not measure time in the development, since its tasks ran in parallel.

\section{Kept Changes}
\label{sec:ri_kept}

We keep three changes, each applied to the model that results from the previous ones. \Cref{tab:ri_chain} compares every step of the chain with the step before it.

\begin{table}[htbp]
\centering
\caption{The kept changes, each applied to the step above it, on the development runs: medians over the 32 runs less the failed ones. Each row is compared with the row above it over the runs that both finished: a value in bold is significantly better than that of the row above it, and a value in red significantly worse; a smaller summary is better. A dash marks the Steady-State ARI of the streams whose penultimate batch holds a single class. Every failure is a summary that exceeded a quarter of the tuning segment.}
\adjustbox{width=\textwidth, totalheight=0.82\textheight, keepaspectratio}{%
\begin{tabular}{l c c c c c c }
\toprule
Model & INS-A & INS-I & INS-R & GAS & PEN & SHU \\
\midrule
\multicolumn{7}{l}{\textit{Overall ARI}} \\
\modelname{MicroTEDAClus+ (Ch.~\ref{chap:ModelImprovement})} & 0.011 & 0.087 & 0.013 & 0.163 & 0.410 & 0.458 \\
\modelname{low-variance} & \textbf{0.142} & \textbf{0.137} & \textbf{0.107} & 0.154 & \textbf{0.505} & 0.417 \\
\modelname{low-variance + no-noise} & 0.153 & 0.137 & \textbf{0.144} & 0.174 & 0.492 & 0.404 \\
\modelname{low-variance + no-noise + link-factor} & 0.140 & \textcolor{red}{0.111} & \textcolor{red}{0.110} & \textbf{0.203} & \textbf{0.577} & 0.469 \\
\midrule
\multicolumn{7}{l}{\textit{Steady-State ARI}} \\
\modelname{MicroTEDAClus+ (Ch.~\ref{chap:ModelImprovement})} & - & 0.104 & - & 0.093 & 0.377 & 0.415 \\
\modelname{low-variance} & - & 0.164 & - & 0.076 & \textbf{0.543} & 0.365 \\
\modelname{low-variance + no-noise} & - & 0.168 & - & 0.100 & 0.512 & 0.308 \\
\modelname{low-variance + no-noise + link-factor} & - & 0.160 & - & \textbf{0.135} & 0.558 & 0.437 \\
\midrule
\multicolumn{7}{l}{\textit{Steady-state summary}} \\
\modelname{MicroTEDAClus+ (Ch.~\ref{chap:ModelImprovement})} & 78 & 93 & 9 & 112 & 71 & 193 \\
\modelname{low-variance} & \textcolor{red}{451} & \textcolor{red}{106} & \textcolor{red}{215} & 122 & \textcolor{red}{263} & 161 \\
\modelname{low-variance + no-noise} & 671 & 107 & \textcolor{red}{709} & 122 & 259 & 129 \\
\modelname{low-variance + no-noise + link-factor} & 964 & 294 & \textbf{309} & 128 & 286 & \textcolor{red}{383} \\
\midrule
\multicolumn{7}{l}{\textit{Failed runs}} \\
\modelname{MicroTEDAClus+ (Ch.~\ref{chap:ModelImprovement})} & 0 & 0 & 0 & 0 & 0 & 0 \\
\modelname{low-variance} & 1 & 1 & 0 & 0 & 1 & 0 \\
\modelname{low-variance + no-noise} & 4 & 1 & 0 & 0 & 4 & 0 \\
\modelname{low-variance + no-noise + link-factor} & 1 & 0 & 0 & 0 & 2 & 1 \\
\bottomrule
\end{tabular}
}
\label{tab:ri_chain}
\end{table}


\subsection{Variance Limit}
\label{subsec:ri_low_variance}

\modelname{low-variance} searches the variance limit from a tenth of the lower end of the range of the relative variance, \cref{tab:real_rules}. With the range of \cref{chap:real_data}, the selected limit lies at the lower end of its range in most runs of \InsectsAbrupt, \InsectsIncremental and \Pendigits, \cref{tab:ri_variance}, so the search would select smaller limits if it could. With the extended range, it lies at the lower end in no run.

\begin{table}[htbp]
\centering
\caption{Percentage of the development runs in which the selected variance limit lies within a twentieth of the width of its range from its lower end, in the logarithmic scale of the search.}
\adjustbox{width=\textwidth, totalheight=0.82\textheight, keepaspectratio}{%
\begin{tabular}{l c c c c c c }
\toprule
Model & INS-A & INS-I & INS-R & GAS & PEN & SHU \\
\midrule
\multicolumn{7}{l}{\textit{Variance limit at the lower end of its range}} \\
\modelname{MicroTEDAClus+ (Ch.~3)} & 75 & 84 & 34 & 0 & 81 & 0 \\
\modelname{low-variance} & 0 & 0 & 0 & 0 & 0 & 0 \\
\bottomrule
\end{tabular}
}
\label{tab:ri_variance}
\end{table}


The change widens the search of \MicroTEDAClusPlus and of no other model. The ranges of \cref{chap:real_data} come from the tuning segment by rule, so a wider range does not suit one scale only. In that chapter, we extended the lower end of the distances of the benchmark models to a tenth of the median distance from a sample to its nearest neighbour, \cref{tab:real_rules}, and left the lower end of the relative variance at the square of that distance.

\modelname{low-variance} improves \MicroTEDAClusPlus of \cref{chap:ModelImprovement} in the Overall ARI of the three INSECTS datasets and of \Pendigits, and in the Steady-State ARI of \Pendigits, \cref{tab:ri_chain}. Its steady-state summary is larger where it improves.

\subsection{Noise Label}
\label{subsec:ri_no_noise}

\modelname{no-noise} removes the noise label of \MicroTEDAClusPlus, \cref{eq:mi_final_prediction}: the model assigns every sample to the macrocluster of its nearest active microcluster, and the noise distance is no longer a hyperparameter. The classes of the real datasets cover every sample, \cref{sec:real_datasets}, so a sample labelled as noise is always an error of the partition, and the previous version of \MicroTEDAClusPlus labels nearly half of the samples of \InsectsAbrupt as noise, \cref{tab:ri_diagnosis}.

\modelname{low-variance + no-noise} improves \modelname{low-variance} in the Overall ARI of \InsectsReoccurring, and loses in no dataset, \cref{tab:ri_chain}. It fails in more runs, and its summary is larger in \InsectsReoccurring.

\subsection{Link Factor}
\label{subsec:ri_link}

\modelname{link-factor} is the candidate of \cref{subsec:mi_improvement_candidates} that searches the factor $\ell$ of the intersection condition, \cref{eq:microcluster_intersection_condition}, in $[0.5, 4]$. Applied to \modelname{low-variance}, it improves it with a single gain, while its median Overall ARI is lower in four of the six datasets, \cref{tab:ri_candidates}. Since we kept \modelname{no-noise} against the same control, we evaluate the two changes together: we keep \modelname{link-factor} if we keep \modelname{low-variance + no-noise + link-factor} against \modelname{low-variance + no-noise}.

The combination is mixed, \cref{tab:ri_chain}: it gains in the Overall ARI of \GasDrift and \Pendigits and in the Steady-State ARI of \GasDrift, and loses in the Overall ARI of \InsectsIncremental and \InsectsReoccurring. The search selects factors close to the fixed factor of two in the INSECTS datasets and different factors in the others, \cref{tab:ri_link}.

\begin{table}[htbp]
\centering
\caption{Link factor selected on the development runs: median and quartile coefficient of dispersion.}
\adjustbox{width=\textwidth, totalheight=0.82\textheight, keepaspectratio}{%
\begin{tabular}{l c c c c c c }
\toprule
 & INS-A & INS-I & INS-R & GAS & PEN & SHU \\
\midrule
\multicolumn{7}{l}{\textit{\modelname{low-variance + no-noise + link-factor}}} \\
Link factor & 2.40 [0.20] & 2.14 [0.20] & 2.04 [0.21] & 2.84 [0.17] & 1.17 [0.09] & 3.41 [0.15] \\
\bottomrule
\end{tabular}
}
\label{tab:ri_link}
\end{table}


Where the selected factor is that of the control, the model is nearly the control, found by a search with one more hyperparameter and the same number of trials, which we already found to cost quality in \cref{subsec:mi_m_function}. We attribute the losses to the search and not to the change, and the gains outnumber them, so we keep the combination.

\section{Rejected Candidates}
\label{sec:ri_rejected}

\Cref{tab:ri_candidates} gives the outcome of every other candidate against its control. Its last rows are \modelname{link-factor} applied alone to \modelname{low-variance}, before the combination test of \cref{subsec:ri_link}, and \modelname{wide-link} against its control and against the last step of the chain.

\begin{table}[htbp]
\centering
\small
\caption[The candidates outside the chain of kept changes]{The candidates outside the chain of kept changes, each against its control on the development runs: outcome of the decision rule, significant gains and losses, datasets where its median Overall ARI is lower, and datasets where its steady-state summary is significantly larger or smaller. The covariance was evaluated on eight runs of three datasets.}
\adjustbox{max width=\textwidth, max totalheight=0.82\textheight}{%
\begin{tabular}{l l c c c c c c}
\toprule
Candidate & Control & Outcome & Gains & Losses & Lower & Larger & Smaller \\
\midrule
\variant{m-fixed} & \variant{low-variance} & regresses & 0 & 7 & 6 & 0 & 6 \\
\variant{m-constant} & \variant{low-variance} & regresses & 0 & 7 & 6 & 0 & 6 \\
\variant{m-sigmoid} & \variant{low-variance} & regresses & 0 & 5 & 5 & 0 & 0 \\
\variant{m-reparam} & \variant{low-variance} & regresses & 0 & 2 & 3 & 0 & 0 \\
\variant{m-rate} & \variant{low-variance} & regresses & 0 & 2 & 5 & 0 & 0 \\
\variant{chi-threshold} & \variant{low-variance} & improves & 1 & 0 & 5 & 0 & 0 \\
\variant{diagonal-variance} & \variant{low-variance} & regresses & 0 & 5 & 5 & 0 & 2 \\
\variant{covariance} & \variant{low-variance} & regresses & 0 & 2 & 2 & 0 & 1 \\
\variant{volume-density} & \variant{low-variance} & mixed & 1 & 1 & 4 & 1 & 0 \\
\variant{log-density} & \variant{low-variance} & regresses & 0 & 6 & 6 & 0 & 2 \\
\variant{diagonal-variance + log-density} & \variant{low-variance} & not run & - & - & - & - & - \\
\variant{no-activation} & \variant{low-variance} & mixed & 2 & 6 & 5 & 4 & 0 \\
\variant{mature-mean} & \variant{low-variance} & indistinguishable & 0 & 0 & 4 & 1 & 0 \\
\variant{recency} & \variant{low-variance} & mixed & 2 & 3 & 5 & 2 & 0 \\
\variant{dimensionality-penalty} & \variant{low-variance} & mixed & 2 & 4 & 5 & 0 & 6 \\
\variant{norm-typicality} & \variant{low-variance} & regresses & 0 & 10 & 6 & 0 & 2 \\
\variant{local-weights} & \variant{low-variance} & regresses & 0 & 9 & 6 & 0 & 0 \\
\variant{link-factor} & \variant{low-variance} & improves & 1 & 0 & 4 & 1 & 1 \\
\variant{wide-link} & \variant{low-variance + no-noise} & mixed & 1 & 1 & 3 & 1 & 3 \\
\variant{wide-link} & \variant{low-variance + no-noise + link-factor} & improves & 2 & 0 & 5 & 0 & 3 \\
\bottomrule
\end{tabular}
}
\label{tab:ri_candidates}
\end{table}


\paragraph{Sensitivity function.} \modelname{m-constant} and \modelname{m-reparam} are the candidates of \cref{subsec:mi_m_function}, and three more change the function $m$ of \cref{eq:m_function}: \modelname{m-fixed} sets $m = 3$ from the third sample on, \modelname{m-sigmoid} searches the two constants of the sigmoid, and \modelname{m-rate} keeps the asymptote of three and searches how fast $m$ approaches it. All five regress, and we retain the original function, as in \cref{chap:ModelImprovement}.

\paragraph{Geometry.} \modelname{diagonal-variance} gives a microcluster one variance per feature, and \modelname{covariance} is the candidate of \cref{subsec:mi_covariance} with the directional link. Both regress. We evaluated \modelname{covariance} first on eight runs of \InsectsIncremental, \Pendigits and \Shuttle, three datasets with fewer features than \GasDrift, because each of its microclusters carries a matrix of the square of the dimension and a full development would have taken several hours, most of them in \GasDrift. Its probe already lost in two cells, and \modelname{diagonal-variance}, which also standardizes the distance feature by feature, had regressed in five, so we did not run the full development. We retain the scalar variance.

\paragraph{Density.} \modelname{volume-density} is the density of \cref{sec:mi_corrections}, and \modelname{log-density} is the logarithm of the number of samples per volume of the microcluster. \modelname{volume-density} is mixed with one gain and one loss, which the rule does not keep, and \modelname{log-density} regresses. We did not run \modelname{diagonal-variance + log-density}, since both of its changes regress alone and the pair would have to undo both regressions to be kept.

\paragraph{Activation and pruning.} Two candidates of \cref{subsec:mi_improvement_candidates}, \modelname{no-activation} and \modelname{recency}, are mixed, with more losses than gains. \modelname{mature-mean}, the correction of \cref{sec:mi_corrections}, is indistinguishable from its control and analytically preferred, but its summary is larger in one dataset, so we do not keep it.

\paragraph{Membership.} \modelname{norm-typicality} and \modelname{local-weights} are the memberships of \cref{sec:mi_corrections}. Both regress, and we retain the nearest membership.

\paragraph{Dimensionality penalty.} \modelname{dimensionality-penalty} restores the factor $2/d$ that we removed from \MicroTEDAClusPlus in \cref{chap:ModelImprovement}. It is mixed, with more losses than gains, and its summary is smaller than that of its control in every dataset.

\paragraph{Departures from the rule.} We did not keep two candidates that the rule would keep. \modelname{chi-threshold}, the candidate of \cref{subsec:mi_improvement_candidates} that tests the outliers against the dispersion of a Gaussian cluster in $d$ dimensions, improves \modelname{low-variance} with a single gain, \cref{tab:ri_candidates}, while its median Overall ARI is lower in five of the six datasets. \modelname{wide-link}, the link factor searched in $[0.05, 4]$ on a logarithmic scale, is mixed against its control, \modelname{low-variance + no-noise}, with one gain and one loss, but in the direct comparison planned between it and \modelname{low-variance + no-noise + link-factor} it improves the latter with two gains, while its median Overall ARI is lower in five of the six datasets. In both cases we judged a single favourable dataset to weigh less than a consistent lean in the others, and kept the model without the candidate. We also compared the medians of the two link factors with those of the benchmark models in \cref{tab:real_quality}, outside the rule, and that comparison favoured the narrower range. We kept the link factor itself only provisionally, although it improved its control, because its median Overall ARI was also lower in most datasets; what distinguished it from \modelname{chi-threshold} was the size of its gain, which the rule does not consider. The combination test of \cref{subsec:ri_link} then kept it by the rule. Finally, we did not evaluate \modelname{covariance} and \modelname{diagonal-variance + log-density} in full, for the reasons given with them above.

\section{Confirmation}
\label{sec:ri_confirmation}

\begin{table}[htbp]
\centering
\small
\caption[The kept changes, each applied to the step above it, on the runs of the real-data experiment]{The kept changes, each applied to the step above it, on the runs of the real-data experiment: medians over the 32 runs less the failed ones, each row compared with the row above it over the runs that both finished. A dash marks the Steady-State ARI of a stream whose penultimate batch holds a single class; every failure is a summary over its limit.}
\adjustbox{max width=\textwidth, max totalheight=0.82\textheight}{%
\begin{tabular}{l c c c c c c c c }
\toprule
Model & INS-A & INS-I & INS-R & GAS & KDD & COV & PEN & SHU \\
\midrule
\multicolumn{9}{l}{\textit{Overall ARI}} \\
\modelname{MicroTEDAClus+ (Ch.~\ref{chap:ModelImprovement})} & 0.017 & 0.088 & 0.012 & 0.162 & 0.669 & 0.000 & 0.414 & 0.444 \\
\variant{low-variance} & \textbf{0.138} & \textbf{0.116} & \textbf{0.124} & 0.164 & 0.671 & 0.000 & \textbf{0.530} & 0.433 \\
\variant{low-variance + no-noise} & 0.152 & 0.133 & \textbf{0.144} & 0.170 & 0.697 & 0.052 & 0.500 & 0.468 \\
\variant{low-variance + no-noise + link-factor} & 0.142 & 0.128 & \textcolor{red}{0.106} & 0.192 & 0.689 & \textbf{0.059} & \textbf{0.581} & 0.500 \\
\midrule
\multicolumn{9}{l}{\textit{Steady-State ARI}} \\
\modelname{MicroTEDAClus+ (Ch.~\ref{chap:ModelImprovement})} & - & 0.103 & - & 0.097 & 0.644 & - & 0.356 & 0.399 \\
\variant{low-variance} & - & 0.161 & - & 0.092 & 0.657 & - & \textbf{0.541} & 0.405 \\
\variant{low-variance + no-noise} & - & 0.164 & - & 0.129 & 0.665 & - & 0.515 & 0.422 \\
\variant{low-variance + no-noise + link-factor} & - & 0.153 & - & 0.142 & 0.637 & - & 0.576 & 0.460 \\
\midrule
\multicolumn{9}{l}{\textit{Steady-state summary}} \\
\modelname{MicroTEDAClus+ (Ch.~\ref{chap:ModelImprovement})} & 81 & 94 & 9 & 100 & 12 & 4 & 74 & 182 \\
\variant{low-variance} & \textcolor{red}{326} & \textcolor{red}{102} & \textcolor{red}{748} & 117 & 12 & 5 & \textcolor{red}{370} & 136 \\
\variant{low-variance + no-noise} & 1018 & 108 & 773 & 112 & 16 & 2138 & 262 & 182 \\
\variant{low-variance + no-noise + link-factor} & 1002 & 304 & \textbf{376} & 83 & 17 & 722 & 226 & 329 \\
\midrule
\multicolumn{9}{l}{\textit{Failed runs}} \\
\modelname{MicroTEDAClus+ (Ch.~\ref{chap:ModelImprovement})} & 0 & 0 & 0 & 1 & 0 & 0 & 0 & 0 \\
\variant{low-variance} & 2 & 0 & 0 & 0 & 0 & 14 & 3 & 0 \\
\variant{low-variance + no-noise} & 2 & 0 & 0 & 0 & 0 & 17 & 3 & 0 \\
\variant{low-variance + no-noise + link-factor} & 0 & 0 & 2 & 0 & 0 & 8 & 1 & 1 \\
\bottomrule
\end{tabular}
}
\label{tab:ri_confirmation}
\end{table}


\Cref{tab:ri_confirmation} repeats the chain on the runs 0 to 31 of the eight datasets. \modelname{low-variance} improves the previous \MicroTEDAClusPlus version in the Overall ARI of the three INSECTS datasets and of \Pendigits and in the Steady-State ARI of \Pendigits, the same gains as in the development, \cref{tab:ri_chain}, and its summary is again larger where it improves. The removal of the noise label improves \modelname{low-variance} in \InsectsReoccurring, as in the development. The link factor gains in \Pendigits and loses in \InsectsReoccurring, as in the development, and also gains in \Covertype; its gain in \GasDrift and its loss in \InsectsIncremental are not significant on these runs.

No step loses where the development found it to gain, and no step loses in the two datasets that took no part in the development. In \KDDCup no step changes the model significantly. In \Covertype the steps fail in many runs, all by the limit on the summary of \cref{sec:real_protocol}, \cref{tab:ri_confirmation}.

\section{The Final Version of \MicroTEDAClusPlus}
\label{sec:ri_final}

The final version of \MicroTEDAClusPlus differs from the version of \cref{chap:ModelImprovement} in the three points of \cref{tab:ri_differences}. Its hyperparameters are the relative variance limit $r$, the decay and the link factor $\ell$.

\begin{table}[htbp]
\centering
\caption{Differences between the final version of \MicroTEDAClusPlus and the version of \cref{chap:ModelImprovement}.}
\begin{tabular}{l p{0.36\textwidth} p{0.36\textwidth}}
\toprule
 & Version of \cref{chap:ModelImprovement} & Final version \\
\midrule
Noise & a sample farther than $n_\sigma$ standard deviations from every active microcluster is noise & every sample is assigned to the macrocluster of its nearest active microcluster \\
Intersection & $\lVert \mu_i - \mu_j \rVert \le 2\,(\sigma_i + \sigma_j)$ & $\lVert \mu_i - \mu_j \rVert \le \ell\,(\sigma_i + \sigma_j)$, with $\ell$ searched in $[0.5, 4]$ \\
Variance limit & searched from the square of the median distance from a sample to its nearest neighbour, as a fraction of the total variance & searched from a tenth of that lower end \\
\bottomrule
\end{tabular}
\label{tab:ri_differences}
\end{table}

We tune and evaluate the final version by the protocol of \cref{chap:real_data} on its runs 0 to 31, as every model of that chapter, and measure its time and memory, and those of the benchmark models, by the procedure of \cref{sec:mod_measurements} on the runs 0 to 9. Its search uses seeds of its own, as that of every model of \cref{chap:real_data}, where the confirmation used those of \MicroTEDAClusPlus of \cref{chap:ModelImprovement} to pair the steps, so its medians differ somewhat from those of the last step of \cref{tab:ri_confirmation}.

\begin{table}[htbp]
\centering
\caption{The partitions of the two versions of \MicroTEDAClusPlus{} on the real datasets: medians over the batches of the evaluation segment, then medians over the runs, at least 32 runs, except 27 in COV, 31 in GAS, 31 in INS-A, 31 in INS-R, 30 in PEN. Homogeneity and completeness compare the clusters predicted for a batch with its classes, over the samples not labelled as noise.}
\adjustbox{width=\textwidth, totalheight=0.82\textheight, keepaspectratio}{%
\begin{tabular}{l c c c c c c c c }
\toprule
Model & INS-A & INS-I & INS-R & GAS & KDD & COV & PEN & SHU \\
\midrule
\multicolumn{9}{l}{\textit{Clusters per batch}} \\
Classes & 5.0 & 6.0 & 5.0 & 6.0 & 4.0 & 2.0 & 10.0 & 3.0 \\
\modelname{MicroTEDAClus+ (Ch.~3)} & 3.0 & 7.0 & 1.0 & 9.0 & 1.0 & 1.0 & 8.0 & 8.5 \\
\modelname{MicroTEDAClus+} & 27.0 & 14.0 & 13.0 & 9.0 & 2.0 & 19.0 & 18.0 & 6.0 \\
\midrule
\multicolumn{9}{l}{\textit{Share of the samples labelled as noise}} \\
\modelname{MicroTEDAClus+ (Ch.~3)} & 0.45 & 0.10 & 0.15 & 0.03 & 0.44 & 0.00 & 0.10 & 0.02 \\
\modelname{MicroTEDAClus+} & 0.00 & 0.00 & 0.00 & 0.00 & 0.00 & 0.00 & 0.00 & 0.00 \\
\midrule
\multicolumn{9}{l}{\textit{Homogeneity}} \\
\modelname{MicroTEDAClus+ (Ch.~3)} & 0.09 & 0.27 & 0.00 & 0.40 & 1.00 & 0.00 & 0.65 & 0.61 \\
\modelname{MicroTEDAClus+} & 0.60 & 0.42 & 0.42 & 0.38 & 0.59 & 0.71 & 0.83 & 0.57 \\
\midrule
\multicolumn{9}{l}{\textit{Completeness}} \\
\modelname{MicroTEDAClus+ (Ch.~3)} & 0.51 & 0.36 & 1.00 & 0.47 & 1.00 & 1.00 & 0.78 & 0.41 \\
\modelname{MicroTEDAClus+} & 0.32 & 0.35 & 0.30 & 0.41 & 0.86 & 0.16 & 0.71 & 0.48 \\
\bottomrule
\end{tabular}
}
\label{tab:ri_partitions}
\end{table}


\Cref{tab:ri_partitions} compares the partitions that the two versions predict. The final version labels no sample as noise, and in every dataset but \KDDCup it predicts more clusters than there are classes, several times more in \InsectsAbrupt and \Covertype. Its homogeneity is higher than that of the version of \cref{chap:ModelImprovement} in five datasets and its completeness is lower in most: the final version splits the classes over more clusters, where the version of \cref{chap:ModelImprovement} merged them or labelled their samples as noise. Where the version of \cref{chap:ModelImprovement} predicts a single cluster, in \InsectsReoccurring, \KDDCup and \Covertype, its completeness is one by construction.

\section{Comparison with the Benchmark Models}
\label{sec:ri_comparison}

\subsection{Clustering Quality}
\label{subsec:ri_quality}

\begin{table}[htbp]
\centering
\small
\caption[Clustering quality of \modelname{MicroTEDAClus+} and the benchmark models on the real datasets]{Clustering quality of \modelname{MicroTEDAClus+} and the benchmark models on the real datasets: medians over 94 runs, except 90 in INS-A, 93 in INS-I, 59 in INS-R, 93 in GAS, 54 in KDD, 70 in COV, 81 in PEN. A dash marks a value left out: the Steady-State ARI of a stream whose penultimate batch holds a single class, or a model that failed in more than half of the runs of a dataset.}
\adjustbox{max width=\textwidth, max totalheight=0.82\textheight}{%
\begin{tabular}{l c c c c c c c c }
\toprule
Model & INS-A & INS-I & INS-R & GAS & KDD & COV & PEN & SHU \\
\midrule
\multicolumn{9}{l}{\textit{Overall ARI}} \\
\modelname{MicroTEDAClus+} & 0.147 & 0.134 & 0.105 & 0.184 & 0.665 & 0.058 & 0.585 & 0.463 \\
Better/same/worse & 0/2/2 & 1/0/3 & 1/2/2 & 2/1/1 & 1/1/3 & 0/1/4 & 1/4/0 & 5/0/0 \\
\modelname{DenStream} & \textcolor{red}{0.169} & \textcolor{red}{0.174} & \textcolor{red}{0.153} & 0.184 & \textcolor{red}{0.946} & 0.056 & \textbf{0.506} & \textbf{0.390} \\
\modelname{DBSTREAM} & 0.156 & \textbf{0.089} & 0.084 & \textbf{0.151} & \textcolor{red}{0.916} & \textcolor{red}{0.069} & 0.539 & \textbf{0.364} \\
\modelname{BIRCH} & - & - & \textbf{0.062} & - & \textbf{0.066} & \textcolor{red}{0.065} & 0.530 & \textbf{0.000} \\
\modelname{Mini-batch k-means} & 0.173 & \textcolor{red}{0.167} & 0.115 & \textcolor{red}{0.206} & \textcolor{red}{0.909} & \textcolor{red}{0.073} & 0.529 & \textbf{0.428} \\
\modelname{BICO} & \textcolor{red}{0.182} & \textcolor{red}{0.162} & \textcolor{red}{0.124} & \textbf{0.157} & 0.562 & \textcolor{red}{0.069} & 0.544 & \textbf{0.353} \\
\midrule
\multicolumn{9}{l}{\textit{Steady-State ARI}} \\
\modelname{MicroTEDAClus+} & - & 0.165 & - & 0.143 & 0.581 & - & 0.533 & 0.397 \\
Better/same/worse & - & 1/0/3 & - & 1/3/0 & 1/1/3 & - & 0/5/0 & 2/3/0 \\
\modelname{DenStream} & - & \textcolor{red}{0.252} & - & 0.132 & \textcolor{red}{0.941} & - & 0.500 & 0.385 \\
\modelname{DBSTREAM} & - & \textbf{0.108} & - & 0.116 & \textcolor{red}{0.889} & - & 0.501 & 0.326 \\
\modelname{BIRCH} & - & - & - & - & \textbf{0.022} & - & 0.530 & \textbf{0.000} \\
\modelname{Mini-batch k-means} & - & \textcolor{red}{0.222} & - & 0.135 & \textcolor{red}{0.896} & - & 0.532 & 0.378 \\
\modelname{BICO} & - & \textcolor{red}{0.223} & - & \textbf{0.113} & 0.526 & - & 0.551 & \textbf{0.297} \\
\bottomrule
\end{tabular}
}
\label{tab:ri_quality}
\end{table}

\begin{table}[htbp]
\centering
\caption{Number of runs in which \modelname{MicroTEDAClus+} failed on a real dataset. Every failure is a summary that exceeded a quarter of the tuning segment.}
\adjustbox{width=\textwidth, totalheight=0.82\textheight, keepaspectratio}{%
\begin{tabular}{l c c c c c c c c }
\toprule
Model & INS-A & INS-I & INS-R & GAS & KDD & COV & PEN & SHU \\
\midrule
\multicolumn{9}{l}{\textit{Failed runs}} \\
\modelname{MicroTEDAClus+} & 1 & 0 & 1 & 0 & 0 & 5 & 2 & 0 \\
\bottomrule
\end{tabular}
}
\label{tab:ri_failures}
\end{table}


In the Overall ARI, \cref{tab:ri_quality}, the final version of \MicroTEDAClusPlus is indistinguishable from the best model in \Shuttle, where no other model is, and from the worst in \InsectsAbrupt, \GasDrift and \Covertype; in \InsectsAbrupt its median is the lowest of the models compared. In the four other datasets it is in neither group. In the same comparison, \MicroTEDAClusPlus of \cref{chap:ModelImprovement} was indistinguishable from the worst model in six datasets, \cref{tab:real_quality}. In \KDDCup the median of the final version is about that of the version of \cref{chap:ModelImprovement}.

At the penultimate batch, the final version is indistinguishable from the best model in \GasDrift and \Shuttle and from the worst in \InsectsIncremental. At the fifth batch it is indistinguishable from the worst in the three INSECTS datasets and in \KDDCup, so it still starts slowly, and its maximum ARI drop is among the largest in \InsectsIncremental, \Pendigits and \Shuttle. It fails in nine runs, most of them in \Covertype, all by the limit on the summary, \cref{tab:ri_failures}.

\subsection{Cost}
\label{subsec:ri_cost}

\begin{table}[htbp]
\centering
\caption{Execution time and memory of the models on the real datasets, each task measured alone and its memory counted from every allocation of its process: medians over 10 runs, except 9 in INS-A, 6 in INS-R, 4 in KDD, 7 in COV, 9 in PEN. A value in bold is statistically indistinguishable from the best model of its dataset, and a value in red from the worst. A dash marks a model left out of a dataset.}
\adjustbox{width=\textwidth, totalheight=0.82\textheight, keepaspectratio}{%
\begin{tabular}{l c c c c c c c c }
\toprule
Model & INS-A & INS-I & INS-R & GAS & KDD & COV & PEN & SHU \\
\midrule
\multicolumn{9}{l}{\textit{Learning time of the evaluation segment, in milliseconds}} \\
\modelname{MicroTEDAClus+} & \textcolor{red}{2040} & 500 & 3695 & \textcolor{red}{210} & 1970 & 25060 & 50 & \textcolor{red}{1695} \\
\modelname{DenStream} & \textcolor{red}{790} & 1905 & 1165 & \textcolor{red}{785} & 7970 & 3470 & 90 & 220 \\
\modelname{DBSTREAM} & \textcolor{red}{6490} & \textcolor{red}{5670} & 6595 & \textcolor{red}{1115} & 64720 & 26580 & \textcolor{red}{820} & \textcolor{red}{845} \\
\modelname{BIRCH} & - & - & 6510 & - & 12045 & 7470 & 190 & \textcolor{red}{650} \\
\modelname{Mini-batch k-means} & \textbf{170} & 190 & 280 & \textbf{50} & 1910 & 2190 & 30 & 190 \\
\modelname{BICO} & \textbf{120} & \textbf{135} & 200 & 240 & 1215 & 590 & \textbf{10} & \textbf{40} \\
\midrule
\multicolumn{9}{l}{\textit{Prediction time of the evaluation segment, in milliseconds}} \\
\modelname{MicroTEDAClus+} & \textbf{60} & \textbf{30} & 70 & \textbf{10} & 225 & 1290 & \textbf{0} & \textbf{55} \\
\modelname{DenStream} & 300 & 410 & 420 & \textcolor{red}{410} & 2945 & 1460 & 50 & \textbf{70} \\
\modelname{DBSTREAM} & \textcolor{red}{870} & \textcolor{red}{425} & 1505 & \textcolor{red}{335} & 3525 & 5470 & 150 & 145 \\
\modelname{BIRCH} & - & - & 540 & - & 2340 & 2300 & 40 & 190 \\
\modelname{Mini-batch k-means} & \textbf{70} & 80 & 115 & \textbf{20} & 790 & 930 & \textbf{10} & \textbf{80} \\
\modelname{BICO} & \textcolor{red}{710} & \textcolor{red}{770} & 1040 & 255 & 6900 & 8210 & \textcolor{red}{170} & \textcolor{red}{505} \\
\midrule
\multicolumn{9}{l}{\textit{Memory held at the end of the stream, in KiB}} \\
\modelname{MicroTEDAClus+} & \textcolor{red}{191} & 124 & 268 & 449 & 56 & 186 & 34 & \textcolor{red}{92} \\
\modelname{DenStream} & 65 & 101 & 60 & 375 & 58 & 16 & 53 & \textbf{7} \\
\modelname{DBSTREAM} & \textcolor{red}{599} & \textcolor{red}{446} & 252 & \textcolor{red}{304} & 448 & 103 & \textcolor{red}{382} & 22 \\
\modelname{BIRCH} & - & - & 1587 & - & 478 & 44 & \textcolor{red}{169} & \textcolor{red}{44} \\
\modelname{Mini-batch k-means} & \textbf{8} & \textbf{8} & 7 & \textbf{22} & 8 & 8 & \textbf{9} & \textbf{5} \\
\modelname{BICO} & \textcolor{red}{188} & \textcolor{red}{184} & 176 & \textcolor{red}{676} & 331 & 43 & 173 & \textcolor{red}{52} \\
\midrule
\multicolumn{9}{l}{\textit{Peak memory during the evaluation, in KiB}} \\
\modelname{MicroTEDAClus+} & \textcolor{red}{3700} & \textcolor{red}{578} & 4464 & \textcolor{red}{829} & 244 & 12097 & \textcolor{red}{332} & \textcolor{red}{5613} \\
\modelname{DenStream} & \textcolor{red}{630} & \textcolor{red}{2273} & 528 & \textcolor{red}{4474} & 1576 & 142 & \textcolor{red}{364} & 142 \\
\modelname{DBSTREAM} & \textcolor{red}{870} & \textcolor{red}{714} & 480 & 1355 & 700 & 324 & \textcolor{red}{509} & 94 \\
\modelname{BIRCH} & - & - & 10104 & - & 1004 & 97 & \textcolor{red}{377} & 106 \\
\modelname{Mini-batch k-means} & \textbf{34} & \textbf{26} & 24 & \textbf{67} & 22 & 32 & \textbf{45} & \textbf{17} \\
\modelname{BICO} & 397 & \textcolor{red}{420} & 422 & \textcolor{red}{1293} & 554 & 198 & \textcolor{red}{320} & 166 \\
\bottomrule
\end{tabular}
}
\label{tab:ri_cost}
\end{table}


The final version is indistinguishable from the fastest model in the prediction stage in five datasets, and from the slowest in none, \cref{tab:ri_cost}. In the learning stage it is indistinguishable from the slowest in \InsectsAbrupt, \GasDrift and \Shuttle and from the fastest in none; the models indistinguishable from the fastest in that stage are \MiniBatchKMeans and \BICO. The memory it holds at the end of the stream is among the largest in \InsectsAbrupt and \Shuttle, and its peak during the evaluation is among the highest in five datasets. \MiniBatchKMeans holds the least memory in five datasets. These comparisons rest on ten runs, and on fewer where a model failed.

\subsection{Remaining Limits}
\label{subsec:ri_limits}

Three limits of the final version remain. It splits the classes of the INSECTS datasets over many clusters, \cref{tab:ri_partitions}, while its search selects there a link factor close to the fixed factor of the version of \cref{chap:ModelImprovement}, \cref{tab:ri_link}, although a larger factor would join more microclusters. It starts slowly, \cref{tab:ri_quality}. And its summary exceeds the limit in some runs of \Covertype, \cref{tab:ri_failures}, and the peak of its memory is among the highest, \cref{tab:ri_cost}.

We evaluated the final version on the real datasets only, and not on the synthetic datasets of \cref{chap:modernization}. Its removal of the noise label suits data whose classes cover every sample, while on the synthetic datasets with noise we found the noise label to improve the model, \cref{sec:mi_improvements}.

\section{Summary}
\label{sec:ri_summary}

On the real datasets of \cref{chap:real_data}, \MicroTEDAClusPlus of \cref{chap:ModelImprovement} is indistinguishable from the worst model in most datasets. On them, most of its microclusters never mature in the INSECTS datasets and in \Pendigits, and it labels as noise samples of classes that cover the whole stream, \cref{tab:ri_diagnosis}. Of the candidates that we evaluated on six of the datasets, we kept three: the variance limit searched one decade lower, the removal of the noise label and a searched link factor, \cref{tab:ri_chain}. The alternatives to the function $m$, to the isotropic variance, to the density, to the activation and to the membership did not improve the model, \cref{tab:ri_candidates}, as in \cref{chap:ModelImprovement}. The chain held on new runs and on the two datasets left out of the development, \cref{tab:ri_confirmation}.

Tuned and evaluated as every model of \cref{chap:real_data}, the final version of \MicroTEDAClusPlus is indistinguishable from the worst model in three datasets instead of six, and from the best in \Shuttle, \cref{tab:ri_quality}. It still starts slowly and splits the classes of the INSECTS datasets over many clusters, \cref{tab:ri_partitions}. It is among the fastest models in the prediction stage, and among those of the highest peak of memory, \cref{tab:ri_cost}.
